\typeout{NT FILE introduction.tex}%

\chapter{Introduction}
\label{cha:introduction}

\section{Motivation}

The increasing frequency and severity of wildfires, aggravated by climate change, 
pose a growing threat to ecosystems, properties, and human lives, particularly in 
wildland-urban interface areas \cite{unep2022wildfires}. In Portugal, wildfires have burned 
over 3 million hectares between 1980 and 2018, roughly one-third of the country's area,
with a concerning rise in fire events in recent years \cite{jrc2020wildfire}. For instance, 
in 2022 alone, more than 110,000 hectares were scorched across 10,449 fires \cite{mdpi2023fire}, 
and the 2024 fire season saw over 135,000 hectares burned, resulting in nine fatalities 
and the evacuation of thousands \cite{wiki2024ptfires}.

In response, Portugal has implemented strict vegetation management regulations 
to reduce fire hazards around residential areas. Decree-Law 124/2006 \cite{portugal2006}, 
updated by Decree-Law 82/2021 \cite{portugal2021}\cite{houserefuge2021}\cite{houserefuge2024}, 
mandates a 50-meter defensible space around structures, requiring vegetation reduction and 
clearance of at least 5 meters between trees and buildings, with a minimum of 4 meters 
between tree canopies, summarized in Figure \ref{fig:veg-man}.

\begin{figure}[ht]
    \centering
    \includegraphics[width=1\textwidth]{../Figures/Veg-management.png}
    \caption{Summary of legal requirements for vegetation management around buildings \cite{icnf1}\cite{icnf2}\cite{icnf3}.}
    \label{fig:veg-man}
\end{figure}

In practice, housing fire risk assessments often depend on extensive manual fieldwork, making them resource-intensive and inconsistent.
While some recent studies have explored the use of satellite imagery and machine learning for wildfire 
prediction, the majority focus on a broad landscape-level. 
There is a notable lack of approaches targeting house-level risk classification, 
particularly those incorporating high-resolution aerial orthophotos and vegetation-specific spectral indices.

We found no prior work targeting this specific scale with use of multi-spectral aerial data. This research aims to fill that gap by 
developing a novel classification framework that leverages multi-spectral aerial imagery and ensemble 
learning to automate the assessment of vegetation-related fire risk around isolated residential structures.


\section{Background}

To improve fire risk assessment, systematic efforts have been undertaken, such as 
the national survey by ADAI, which analyzed approximately 2,000 isolated houses to 
evaluate vegetation-related fire risk \cite{almeida2020analise}. Houses were selected 
based on areas burned between 2016 and 2020, using data from the ICNF (Instituto de 
Conservação da Natureza e das Florestas), ensuring a geographically balanced sample.

\begin{figure}[ht]
    \centering
    \includegraphics[width=0.9\textwidth]{../Figures/FireRiskClassificationTable.png}
    \caption{Housing fire risk classification based on vegetation characteristics, adapted from \cite{almeida2020analise}.}
    \label{fig:fire-risk-table}
\end{figure}


The classification system defines five fire risk levels: low, moderate, high, very high, 
and extreme. These levels are determined based on both the type of vegetation present 
and its proximity to the residential structure. Only fuel types with sufficient density 
and continuity to sustain and propagate fire were considered; irrigated lawns and managed 
ornamental vegetation were excluded.

The assessment focused on vegetation located within a 50-meter radius of each house, using 
imagery captured between 2018 and 2020. Natural fuels were categorized into two main groups: 
spontaneous vegetation (wildland growth) and non-spontaneous vegetation (human-maintained areas, such as 
gardens or agricultural plots). These were further classified into three subtypes: arboreal, shrub, 
and herbaceous.

\section{Objectives}

The goal of this dissertation is to develop a deep learning-based system for automatically classifying fire risk levels associated 
with vegetation near isolated houses, using high-resolution aerial imagery. 
This system will follow the housing fire risk classification framework established by \gls{adai}, focusing on 
vegetation conditions within a 50-meter buffer around each structure.

To support this objective, the study will make use of orthophotos provided by \gls{dgt} and ground-truth labels 
from \gls{adai}, resulting in a labeled dataset that reflects real-world fire risk conditions with house-level detail.

This work introduces several key innovations:
\begin{itemize}
    \item It explores the integration of multi-spectral input formats, specifically RGB, IRG, and NDVI, into convolutional neural network (CNN) models for fire risk classification.
    \item It compares the performance of multiple deep architectures, including ResNet-50, DenseNet-161, and EfficientNet-B5, each tested under different input configurations.
    \item It implements a custom ensemble learning strategy that aggregates outputs from different models and spectral inputs to improve classification consistency across diverse vegetation settings.
\end{itemize}

Model performance will be evaluated using metrics such as accuracy, precision, recall, and F1-score, with a focus 
on correctly identifying very high and extreme risk areas. This study aims to contribute a data-driven, scalable 
framework for automated housing-level fire risk assessment, 
designed to assist vegetation monitoring and inform wildfire mitigation planning.


\section{Document Structure}

This document is organized into the following chapters:

\begin{itemize}
    \item \textbf{Chapter \ref{cha:concepts}: Fundamental Concepts} - Introduces foundational concepts related to this research, including supervised learning, deep learning, convolutional neural networks (CNNs), evaluation metrics, optimization techniques, and relevant aerial image formats (RGB, IRG, NDVI).
    
    \item \textbf{Chapter \ref{cha:related-work}: Related Work} - Reviews prior work on remote sensing image classification, with an emphasis on CNN-based methods for scene classification. This chapter also presents the core CNN architectures adopted in this research, ResNet-50, DenseNet-161, and EfficientNet-B5, discussing their structural innovations and relevance to the housing fire risk classification task.
    
    \item \textbf{Chapter \ref{cha:methodology}: Methodology} - Describes the overall methodological pipeline in detail. This includes data collection procedures, ground truth labeling strategy, dataset challenges, preprocessing steps, augmentation techniques, and model training processes. It further elaborates on the design of custom classification layers, training strategies, and ensemble learning methods for improving classification robustness.
    
    \item \textbf{Chapter \ref{cha:results}: Results} - Presents and analyzes the experimental results obtained across different architectures and input configurations (RGB, IRG, NDVI). It includes performance comparisons, ensemble results, and a discussion of misclassification patterns and their implications.
    
    \item \textbf{Chapter \ref{cha:conclusions}: Conclusions} - Summarizes the major findings of the study, highlights its contributions to automated fire risk assessment, and outlines opportunities for future research in the domain of remote sensing and deep learning.
\end{itemize}
